The index of substitution saturation, $\Iss=\Hbar/\HFSS$, decides, before any
tree is built, whether an alignment retains phylogenetic signal. The entropy
enters $\E[\Iss]$ only through the rate matrix's spectrum: the index is a
linear functional of $\langle e^{\mu d}\rangle$, collapsing onto mean pairwise
$p$-distance exactly when each eigenvalue's multiplicity-to-composition-weight
ratio is constant. That holds for every reversible model at uniform composition
and for degenerate spectra, but fails for non-uniform composition with a
non-degenerate spectrum, the regime it was calibrated in. To second order it
adds nothing about realized substitutions beyond the uncorrected proportion of
differing sites; outside the collapse it adds the model's composition. The
critical value fails separately: it carries no rate-heterogeneity coordinate,
and near saturation it is pinned within a few percent of its ceiling --- a
fifteen-fold change in the tree length at which recovery fails moves it by
$1.2\%$. A four-taxon transfer of long branches at fixed patristic total makes
this direct: the index is constant along the path, recovery spans its full
range. Mean $\lambda$ collapses the same way. Real data occupy both regimes:
the saturated codon partition of the index's own data set falls between two
readings of its threshold, a genome-scale avian alignment five to six times
below both. The index is not information-free --- it predicts recovery where
$p$-distance does not --- but what it adds tracks the tree's shape and site
composition, not the substitution count. Passing the test is evidence against
saturation, failing it is weak evidence, neither is evidence of resolution.
